%!TEX TS-program = xelatex
% Bu dosya sjuoz_word2latex.py tarafından otomatik üretildi.
\documentclass[9pt]{extarticle}
\usepackage{multicol}
\raggedcolumns  % sütun boyu doldurulmaz: başlık ve paragraf aralarında büyük boşluk açılmaz
\usepackage[top=0.5in, bottom=0.75in, right=0.75in, left=0.75in, headheight=0.25in, includehead]{geometry}
\usepackage{fontspec}
\usepackage{amsmath}
\usepackage{amssymb}
\usepackage{graphicx}
% Denklem sütuna sığmazsa: önce 2, sonra 3 satıra bölünmüş biçimi; o da sığmazsa küçültülür
\newsavebox{\sjuozeqbox}
\newlength{\sjuozeqmax}
\newcommand{\sjuozeqset}{\setlength{\sjuozeqmax}{\linewidth}\addtolength{\sjuozeqmax}{-2.5em}}
\newcommand{\sjuozeqput}{\ifdim\wd\sjuozeqbox>\sjuozeqmax
  \typeout{SJUOZ-EQ-SCALED: \theequation; \the\sjuozeqmax; \the\wd\sjuozeqbox}% kontrol raporu için
  \resizebox{\sjuozeqmax}{!}{\usebox{\sjuozeqbox}}\else\usebox{\sjuozeqbox}\fi}
\newcommand{\sjuozfit}[3]{\sjuozeqset\sbox{\sjuozeqbox}{$\displaystyle #1$}%
  \ifdim\wd\sjuozeqbox>\sjuozeqmax\sbox{\sjuozeqbox}{$\displaystyle #2$}\fi
  \ifdim\wd\sjuozeqbox>\sjuozeqmax\sbox{\sjuozeqbox}{$\displaystyle #3$}\fi\sjuozeqput}
\newcommand{\sjuozfitone}[1]{\sjuozeqset\sbox{\sjuozeqbox}{$\displaystyle #1$}\sjuozeqput}
\usepackage{xcolor}
\usepackage{array}
\usepackage{tabularx}
\usepackage{xltabular}
\usepackage{pdflscape}
\usepackage{booktabs}
\usepackage{multirow}
\usepackage{caption}
\usepackage{subcaption}
\usepackage{titlesec}
\usepackage{fancyhdr}
\usepackage{qrcode}
\usepackage{xurl}
\usepackage{lastpage}
\usepackage{indentfirst}
\usepackage{eso-pic}
\usepackage{tikz}
\usepackage[colorlinks=true, urlcolor=blue, linkcolor=black, citecolor=black]{hyperref}

\titleformat{\section}{\large\bfseries\centering}{\thesection.}{1.0em}{}
\titleformat{\subsection}{\large\bfseries}{\thesubsection.}{0.5em}{}
\titleformat{\subsubsection}{\large\bfseries}{\thesubsubsection.}{0.5em}{}
\titlespacing*{\section}{0pt}{8pt}{6pt}
\titlespacing*{\subsection}{0pt}{8pt}{4pt}
\titlespacing*{\subsubsection}{0pt}{6pt}{3pt}

\newcommand{\titlesize}{\fontsize{12}{10}}
\renewcommand{\normalsize}{\fontsize{9pt}{9pt}\selectfont}
\renewcommand{\large}{\fontsize{10pt}{10pt}\selectfont}
\renewcommand{\small}{\fontsize{8pt}{8pt}\selectfont}
\DeclareCaptionLabelFormat{sjuozparen}{#1 (#2)}
\captionsetup[table]{labelfont=bf, justification=raggedright, singlelinecheck=false}
\captionsetup[figure]{labelfont=bf, justification=centering, singlelinecheck=false}
\captionsetup[sub]{labelformat=parens, labelfont=bf, font=small, justification=centering}
\definecolor{sjuozblue}{HTML}{C1D1EB}
\definecolor{colorabs}{HTML}{E2EAF4}

% Yazı tipleri: Times New Roman / Garamond yoksa benzerleri kullanılır
\IfFontExistsTF{Times New Roman}{\setmainfont{Times New Roman}}{%
  \IfFontExistsTF{TeX Gyre Termes}{\setmainfont{TeX Gyre Termes}}{%
    \IfFontExistsTF{Liberation Serif}{\setmainfont{Liberation Serif}}{}}}
\IfFontExistsTF{Garamond}{\newfontfamily{\garamond}{Garamond}}{%
  \IfFontExistsTF{EB Garamond}{\newfontfamily{\garamond}{EB Garamond}}{\newcommand{\garamond}{\rmfamily}}}

\newcommand{\sjuozdoi}{https://doi.org/10.25271/sjuoz.2026.14.4.2012}
% İlk sayfadaki bilgi kutusunun sayfa altından yüksekliği ve metin için ayrılan boşluk
\newcommand{\firstfooterbottom}{0.4in}
\newcommand{\firstpagereserve}{1.4in}
% Sadece ilk sayfada metin alanını kısalt; ilk sayfa basıldıktan sonra eski haline döndür
\makeatletter
\newlength{\sjuozfulltextheight}
\AtBeginDocument{%
  \setlength{\sjuozfulltextheight}{\textheight}%
  \global\advance\textheight by -\firstpagereserve
  \global\@colht\textheight \global\@colroom\textheight \global\vsize\textheight
  \AddToHookNext{shipout/after}{\global\textheight\sjuozfulltextheight}}
\makeatother

% Başlık ve altlık (header/footer)
\pagestyle{fancy}
\fancyhf{}
\renewcommand{\headrulewidth}{0pt}
\lhead{\large \textit{Omer \& Ali}}
\rhead{\small{\garamond{\textbf{SJUOZ|VOL14|OCT 2026|Pxxx-xxx}}}}
\lfoot{\raisebox{8pt}{\textit{DOI:}\href{\sjuozdoi}{\textit{\small{\sjuozdoi}}}}}
\cfoot{\raisebox{8pt}{\thepage}}

\fancypagestyle{firstpage}{%
  \fancyhf{}
  \renewcommand{\headrulewidth}{0pt}
  \rhead{\IfFileExists{SJUOZ.png}{\includegraphics[width=1.04in]{SJUOZ.png}}{\textbf{SJUOZ}}}
  \lhead{\colorbox{sjuozblue}{%
      \begin{minipage}[b]{0.83\textwidth}
        \vspace{0.15cm}
        \hspace{-3pt}\ \normalsize \noindent \garamond{Science Journal of University of Zakho}\\
        \normalsize \textit{\garamond{DOI: \sjuozdoi}}
        \vspace{0.2cm}
      \end{minipage}}}
}

% İlk sayfanın altındaki bilgi kutusu (sayfaya mutlak konumla yerleştirilir)
\newcommand{\sjuozfirstfooter}{%
  \begin{minipage}[b]{\textwidth}
      \rule{\textwidth}{1.5pt} \\[0.15cm]
      \begin{minipage}[t]{0.2\textwidth}\hspace{1pt} \small Access this article online \\
        \vspace{0.1cm}
        {\hypersetup{urlcolor=black}\qrcode[height=2.6162cm]{https://sjuoz.uoz.edu.krd/index.php/sjuoz/article/view/2012}}
      \end{minipage}\begin{minipage}[t]{0.4\textwidth}
        \small DOI:\href{\sjuozdoi}{\sjuozdoi} \\ \\[0.5cm]
        Printed ISSN 2663-628X; \\
        Electronic ISSN 2663-6298
      \end{minipage}\begin{minipage}[t]{0.4\textwidth}\small \textbf{Science Journal of University of Zakho} \\
        \textbf{Vol. 14, No. 04, pp. xxx-xxx, October-2026} \\[0.5cm]
        This is an open access article under a CC BY-NC-SA 4.0 license \\
        (\href{https://creativecommons.org/licenses/by-nc-sa/4.0/}{https://creativecommons.org/licenses/by-nc-sa/4.0/})
      \end{minipage} \\[0.15cm]
      \rule{\textwidth}{1.5pt} \\[0.15cm]
      \centering \normalsize \thepage
  \end{minipage}}

% ORCID simgesi: orcid.png varsa o, yoksa çizilmiş yeşil "iD" simgesi
\definecolor{orcidgreen}{HTML}{A6CE39}
\newcommand{\orcidicon}{\IfFileExists{orcid.png}{\includegraphics[width=9pt]{orcid.png}}{%
  \tikz[baseline=-0.7ex]{\fill[orcidgreen] (0,0) circle (4.5pt);
  \node[text=white, font=\fontsize{4.5pt}{4.5pt}\selectfont\bfseries, inner sep=0pt] at (0,0) {iD};}}}
\newcommand{\orcidlink}[1]{\href{#1}{\orcidicon}}

\begin{document}
\setcounter{page}{1}
\thispagestyle{firstpage}
\AddToShipoutPictureFG*{\put(\LenToUnit{0.75in},\LenToUnit{\firstfooterbottom}){\sjuozfirstfooter}}

\noindent \large \textbf{Original Article} \\[-5pt]
\begin{center} \titlesize \textbf{ON THE CONVERGENCE OF A CLASS OF ITERATIVE METHODS} \end{center}
\begin{center} \large \textbf{Dilan Omer \textsuperscript{1,*} and Karwan Ali \textsuperscript{1}}
\end{center}
\noindent
\large \textsuperscript{1}Department of Mathematics, Faculty of Science, University of Zakho, Duhok, Iraq.\\

\begin{center} \textsuperscript{*}\textit{Corresponding author, E-mail:}\href{mailto:dilan.omer@uoz.edu.krd}{dilan.omer@uoz.edu.krd}
\end{center}
\rule{\textwidth}{1.5pt}
\begin{minipage}[t]{0.175\textwidth}
  \hspace{0cm} \vspace{0cm}
\end{minipage} \begin{minipage}[t]{0.8\textwidth}
\large \textbf{ABSTRACT}\end{minipage} \\[-6pt]
\rule{\textwidth}{1.5pt} \\
\begin{minipage}[t]{0.175\textwidth}
  \vspace{1cm} Received \\ xxx \\[0.5cm] Accepted \\ xxx \\[0.5cm] Published \\ xxx \\[0.25cm]
\end{minipage} \colorbox{colorabs}{\begin{minipage}[t]{0.81\textwidth}
  \large{We study long displayed equations in a two-column layout. Every equation must stay inside its column: long sums are broken at plus and minus signs, chains at the equals sign, and unbreakable expressions are scaled down.
  \\[2pt]
  \normalsize \textbf{Keywords:} iterative methods; convergence; equations.
  }
\end{minipage}}
\\[-10pt]

\normalsize
\noindent \rule{\textwidth}{0.75pt}
\begin{multicols}{2}
\section{INTRODUCTION}
The error of the method satisfies the bound in \hyperlink{eq:1}{Equation (1)}, and the full expansion is given in \hyperlink{eq:2}{Equation (2)}.

\hypertarget{eq:1}{}
\begin{equation}
  \sjuozfitone{\left\| x_{k+1}-x^{*} \right\|\leq\left\| x_{k}-x^{*} \right\|}
\end{equation}
\hypertarget{eq:2}{}
\begin{equation}
  \sjuozfit{F\left( x \right)=a_{0}+a_{1}x+a_{2}x^{2}+a_{3}x^{3}+a_{4}x^{4}+a_{5}x^{5}+a_{6}x^{6}+a_{7}x^{7}+a_{8}x^{8}+a_{9}x^{9}+a_{10}x^{10}+a_{11}x^{11}}{\begin{aligned}F\left( x \right) &=a_{0}+a_{1}x+a_{2}x^{2}+a_{3}x^{3}+a_{4}x^{4}+a_{5}x^{5} \\ &\quad +a_{6}x^{6}+a_{7}x^{7}+a_{8}x^{8}+a_{9}x^{9}+a_{10}x^{10}+a_{11}x^{11}\end{aligned}}{\begin{aligned}F\left( x \right) &=a_{0}+a_{1}x+a_{2}x^{2}+a_{3}x^{3} \\ &\quad +a_{4}x^{4}+a_{5}x^{5}+a_{6}x^{6}+a_{7}x^{7} \\ &\quad +a_{8}x^{8}+a_{9}x^{9}+a_{10}x^{10}+a_{11}x^{11}\end{aligned}}
\end{equation}
The loss combines several terms:

\hypertarget{eq:3}{}
\begin{equation}
  \sjuozfit{L\left( \theta\right)=\frac{1}{N}\sum\limits_{i=1}^{N} {\left( y_{i}-{\hat{y}}_{i} \right)}^{2}+\lambda_{1}\sum\limits_{j=1}^{p} \left| \theta_{j} \right|+\lambda_{2}\sum\limits_{j=1}^{p} \theta_{j}^{2}-\gamma\log\left( 1+e^{-\theta^{T}x} \right)}{\begin{aligned}L\left( \theta\right) &=\frac{1}{N}\sum\limits_{i=1}^{N} {\left( y_{i}-{\hat{y}}_{i} \right)}^{2} \\ &\quad +\lambda_{1}\sum\limits_{j=1}^{p} \left| \theta_{j} \right|+\lambda_{2}\sum\limits_{j=1}^{p} \theta_{j}^{2}-\gamma\log\left( 1+e^{-\theta^{T}x} \right)\end{aligned}}{\begin{aligned}L\left( \theta\right) &=\frac{1}{N}\sum\limits_{i=1}^{N} {\left( y_{i}-{\hat{y}}_{i} \right)}^{2} \\ &\quad +\lambda_{1}\sum\limits_{j=1}^{p} \left| \theta_{j} \right| \\ &\quad +\lambda_{2}\sum\limits_{j=1}^{p} \theta_{j}^{2}-\gamma\log\left( 1+e^{-\theta^{T}x} \right)\end{aligned}}
\end{equation}
A chain of equalities:

\hypertarget{eq:4}{}
\begin{equation}
  \sjuozfit{\int_{0}^{1} \left( 3x^{2}+2x+1 \right)\mathit{dx}={\left[ x^{3}+x^{2}+x \right]}_{0}^{1}=\left( 1+1+1 \right)-\left( 0+0+0 \right)=3}{\begin{aligned}\int_{0}^{1} \left( 3x^{2}+2x+1 \right)\mathit{dx} &={\left[ x^{3}+x^{2}+x \right]}_{0}^{1} \\ &=\left( 1+1+1 \right)-\left( 0+0+0 \right)=3\end{aligned}}{\begin{aligned}\int_{0}^{1} \left( 3x^{2}+2x+1 \right)\mathit{dx} &={\left[ x^{3}+x^{2}+x \right]}_{0}^{1} \\ &=\left( 1+1+1 \right)-\left( 0+0+0 \right)=3\end{aligned}}
\end{equation}
An expression that cannot be broken:

\hypertarget{eq:5}{}
\begin{equation}
  \sjuozfitone{P=\frac{\alpha_{1}\beta_{1}\gamma_{1}+\alpha_{2}\beta_{2}\gamma_{2}+\alpha_{3}\beta_{3}\gamma_{3}+\alpha_{4}\beta_{4}\gamma_{4}+\alpha_{5}\beta_{5}\gamma_{5}}{\delta_{1}\epsilon_{1}+\delta_{2}\epsilon_{2}+\delta_{3}\epsilon_{3}+\delta_{4}\epsilon_{4}}}
\end{equation}
A matrix and cases:

\hypertarget{eq:6}{}
\begin{equation}
  \sjuozfitone{A=\left( \begin{matrix}a_{11} & a_{12} & a_{13} & a_{14} & a_{15} & a_{16} & a_{17} & a_{18} \\ a_{21} & a_{22} & a_{23} & a_{24} & a_{25} & a_{26} & a_{27} & a_{28}\end{matrix} \right)}
\end{equation}
\hypertarget{eq:7}{}
\begin{equation}
  \sjuozfitone{f\left( x \right)=\left\{ \begin{matrix}x^{2}+2x+1, & x\geq0 \\ -x, & x<0\end{matrix} \right.}
\end{equation}
Inline mathematics such as $a_{0}+a_{1}x+a_{2}x^{2}+a_{3}x^{3}+a_{4}x^{4}+a_{5}x^{5}+a_{6}x^{6}$ stays in the text line.

\section{CONCLUSIONS}
All equations stay inside the column (\hyperlink{ref:1}{Smith, 2020}).

\section*{REFERENCES}
\begin{list}{}{
    \setlength{\leftmargin}{1.5em}
    \setlength{\itemindent}{-1.5em}
    \setlength{\parsep}{0pt}
    \setlength{\itemsep}{6pt}
}
\item \hypertarget{ref:1}{}Smith, J. (2020). Long equations in narrow columns. \textit{Journal of Typesetting}, 1(1), 1–10.

\end{list}
\end{multicols}

\end{document}
